\section{Conclusion}\label{sec:conclusion}

Using the spread of repeated LLM forecasts as Black--Litterman view uncertainty is attractive because it replaces a convention with a measurement. We re-examined the framework with point-in-time data, cutoff-clean open models and an out-of-sample choice of $\tau$, and made its claim testable. In the US pilot the dispersion is informative about forecast error beyond volatility, but as $\Om$ it lowers the Sharpe ratio relative to the constant and He--Litterman conventions in every specification. Its ranking is roughly right, but its range is about thirty times too wide. With returns as the only input the views themselves extrapolate short-term momentum and have no predictive value, so the earlier performance gains do not survive. The evaluation protocol applies to any proposed source of view uncertainty. Whether sampled LLM dispersion becomes useful once the views carry information is the question this re-examination leaves open. \pending{revise after the multi-market, multi-model results}
